\section{引言与相关研究}
\subsection{研究问题}
神经网络训练通常写成随机优化问题
\begin{equation}\label{eq:problem}
 \min_{\theta} f(\theta)=\E_{\xi}[\ell(\theta;\xi)].
\end{equation}
其中 $\xi$ 表示一个样本或一个小批量。对大模型而言，优化器不仅决定每一步向哪个方向移动，也会改变状态存储、更新计算量以及训练稳定性。一个每步更复杂的算法可能减少达到目标损失所需的步数，也可能因为更新成本较高而增加实际训练时间。因此，比较优化器时至少需要区分迭代效率与时间效率。

SGD 使用当前小批量梯度更新参数，动量 SGD 则把历史梯度纳入更新方向。在此基础上，Adam 用逐坐标一阶、二阶矩估计调整步长\citep{adam2015}；AdamW 进一步把权重衰减从梯度统计中分离\citep{adamw2019}。这类方法可以处理不同形状的参数，但其基本归一化操作以坐标为单位。神经网络中的线性层权重本来是矩阵，将其视为一个长向量，会弱化输入方向、输出方向及奇异值谱所包含的结构。

Muon 采用另一种做法：先计算动量，再对动量矩阵作近似正交化，最后更新参数\citep{jordan2024}。在精确极分解的理想情况下，所有非零奇异值都被替换为 1，而奇异向量保持不变。更新不再由最大奇异值方向单独支配。不过，这并不直接说明算法必然优于 AdamW：方向归一化同时丢弃幅度信息，噪声可能影响小奇异值方向，近似正交化也有计算成本。

本文围绕三个问题展开：第一，Muon 更新在什么意义下是最速下降方向？第二，在动量与随机梯度同时存在时，可以证明什么收敛结论？第三，在受控矩阵实验与完整 ResNet-18 图像实验中，理论界、精确更新和有限步实现分别表现如何？这三个问题对应几何解释、数学证明和实验检验，而非对所有模型作统一的优劣排序。

\subsection{文献脉络与本文定位}
Bernstein 与 Newhouse\citep{bernstein2024} 从不同范数下的最速下降解释优化算法，为理解矩阵更新提供了直接入口。Cutkosky 与 Mehta\citep{cutkosky2020} 分析归一化随机梯度中的动量作用，说明动量跟踪误差可以成为收敛证明的核心变量。矩阵问题还需要处理谱范数与核范数的对偶关系，以及从核范数误差到 Frobenius 范数误差时出现的维度因子。

早期 Muon 收敛分析\citep{lihong2025} 已讨论精确 SVD 版本，但某些中间不等式需要谨慎核查。后续工作\citep{chang2025} 明确讨论了已有证明的问题，并进一步分析不同动量和方差缩减版本。本文不以早期论文中含矩阵逆的推导为证明依据，而是从范数对偶关系重新组织一条可逐步检查的证明。所得阶数属于已有理论框架的课程性推导，不声称为新的收敛率。

在应用方面，Liu 等\citep{liu2025} 研究了 Muon 向大模型训练扩展时的权重衰减与更新尺度问题。有限步 Newton--Schulz 与 Nesterov 修正的直接分析也已有专门研究\citep{peng2026}。这些工作提醒我们：精确极分解、有限多项式迭代和带其他参数组的混合实现并非同一个数学迭代。报告中必须明确区分。

本文的工作包括：给出处理零矩阵和秩亏情形的精确更新定义；证明精确版本与含谱范数近似误差版本的有限步平均梯度界；在独立调参后比较病态非凸矩阵分解与完整 ResNet-18；用真实更新谱和 $K$ 消融检查近似机制及计算开销。原有受控二次问题、随机梯度和语言模型实验保留在附录。所有数值表格都由运行记录生成，代码入口与数据校验值在第\ref{sec:reproduce}节给出。
